This chapter consequently aims to support known principles such as the Arrhenius-Equation with the collection of data on current findings on temperature-time-regimes as well as humidity and sunlight exposure. \par 





\citet{chepkoech} aimed to assess the impact of storage conditions on the retention of retinol and B-vitamins in selected commercial fortified maize flour. Maize flour is a widely consumed staple food in developing countries and has already been fortified with iron, zinc, retinol, and vitamins B$_1$, B$_2$, B$_3$, B$_6$, B$_9$, and B$_{12}$ to improve their micronutrient intakes. However, there are still many cases of micronutrient deficiencies, although fortification programs have been established. Temperature and relative humidity emerged as two of the main factors of micronutrient stability in maize flour. Across all vitamins analyzed (retinol, thiamine, riboflavin, niacin, folate), higher temperature (35 °C) and higher relative humidity (83\% RH) consistently accelerated degradation. Samples stored at 25 °C / 75\% RH retained substantially more micronutrients over the six‑month period. This effect is likely due to an increase in oxidation rates and isomerization, especially for fat-soluble vitamins such as retinol. Although vitamin levels declined over six months, statistical analysis indicated that there is no significant interaction between storage time and condition for retinol, thiamine, niacin and folate, probably because the differences between the two storage environments (temperature/RH difference) were too small to be statistically significant. With the exception of riboflavin, exhibiting a significant interaction which might be explained with its high sensitivity to light. Nevertheless, the study showed a clear hierarchy of stability: Retinol and thiamine were the most unstable vitamins regarding temperature and humidity, whereas niacin and folate remained comparatively stable. The findings indicate that even moderate increases in temperature and humidity accelerate oxidation and isomerization, confirming the Arrhenius principle that reaction rates rise exponentially with temperature.  \par 

%%%%%Sahin et al. 
\citet{li4daylight} examined how vitamin A (retinyl palmitate) retention in fortified sunflower oil changes over 6 months under different light and packaging conditions. Light was found to be the main driver of vitamin A loss and oxidation.  under daylight, especially in transparent bottles, vitamin A decreased dramatically and in some samples reached zero after 6 months. In contrast, vitamin A remained relatively stable in brown (light-protective) bottles regardless of light exposure. The decline in vitamin A was associated with increasing oil oxidation, which can be linked to interactions between oxidative lipid products and vitamin A. Overall, the best stability occurred when samples were stored in the dark, in brown/light-resistant packaging, and with higher retinyl palmitate fortification levels (about 60-90 µg/g). Sunflower oil  that is stored under daylight, particularly in pellucid bottles, tended to have the lowest oxidative stability and showed the greatest loss of retinyl palmitate. These findings are comparable to observations that humidity and heat accelerate micronutrient degradation in dry food matrices.
From a practical perspective, it is useful to keep sunflower oil away from direct light and store it for no longer than a few months. Additionally, it must be consumed relatively promptly in order to minimize oxidative degradation and preserve vitamin A stability. \par 
%%%%%Essilfie et al. AKA li5cocoyam
\citet{li5cocoyam} focused on minerals in cocoyam leaves, an underutilized but nutrient-rich leafy vegetable, commonly consumed in Ghana. The leaves were treated by using different drying methods (solar drying and freeze drying), blanching pretreatment, two storage temperatures (room: 32 °C and refrigeration: 6.7 °C), and two storage durations (0 and 42 days). Overall, the results show that mineral retention was mainly determined by interactions between processing and storage factors, rather than by one single condition. The study examined five different micronutrients- calcium, phosphorus, sodium, zinc and iron- which revealed considerable differences in stability during storage. 
Calcium was highly unstable during storage, especially at 32 °C, and showed major losses regardless of processing method, while sodium remained relatively stable and was mainly affected by processing. Storage temperature and storage duration did not produce significant changes in sodium content which implies that sodium is the most temperature-stable mineral in the study. Phosphorus was the most variable mineral, strongly influenced by the combined effects of drying method, blanching, storage temperature, and storage time. Both phosphorus and calcium levels depend on interactions between processing and storage factors rather than on any single factor alone. Their stability depends heavily on the interaction of drying method, pretreatment, and temperature. Zinc levels were very low (0,001-0,007 g/ 100 g) and influenced mainly by interaction effects between processing and temperature conditions, with only limited time-related interactions, while iron levels were below the limit of quantification. 
The study highlights that combined effects of drying method, pretreatment, and temperature determine mineral retention. This is again consistent with the Arrhenius‑based temperature‑time dependency.

 \par 
%%%%%Attaugwu
The study indicated that refrigerated storage is a more suitable storage condition than ambient storage. Among the micronutrients that were examined, Vitamin E was the only one that remained relatively stable after a period of 30 days. The Vitamins B$_1$, B$_2$ and B$_3$ displayed reduced micronutrient retention which is aligning with known mechanisms of micronutrient degradation.



%%%\par PASTA FUN WITH YOUR MUM \par 

Pasta was enriched with spinach and carrot powder and millet flour. Retention rates of iron and Vitamin A were observed on storage time which was carried out monthly for four months, as well as storage temperature with relative humidity (RH) of 25 °C; 60\% RH and 40 °C; 80\% RH and consequently on two packaging materials - polypropylene and aluminium laminates.
Foil packed pasta fortified with carrot powder is having greater Vitamin A retention than polypropylene presumably due to sunlight exposure. Vitamin A retention was significantly influenced by packaging, light exposure, temperature and storage time. Whereas iron was significantly influenced by storage time. The data also suggested that iron retention was highest within laminated packaging and ambient storage temperature.

